\subsection{Modalidad de Investigación}

La modalidad de la investigación define el camino y las características del estudio para resolver el problema identificado en el centro odontológico Bellidental. Para garantizar que los resultados sean precisos y útiles, se establecieron tres características metodológicas principales. El proyecto se clasifica como cuantitativo por su forma de manejar los datos, aplicativo por su nivel de intervención tecnológica y preexperimental por su diseño de evaluación.

El enfoque del estudio es cuantitativo porque se basa en la recolección y el análisis numérico de los datos. Esta característica es fundamental para cumplir el tercer objetivo específico del proyecto. Al utilizar este enfoque, es posible medir con números exactos el tiempo que el personal dedica a la gestión administrativa por cada paciente. Comparar estas cifras permite demostrar de manera objetiva y matemática si la integración del flujo de información logra reducir realmente la carga de trabajo.

El alcance de la investigación es aplicativo debido a que no se limita a describir un problema teórico, sino que busca solucionarlo mediante una intervención directa en la realidad de la clínica. Esto se relaciona estrechamente con el primer y segundo objetivo del proyecto. A partir de identificar los cuellos de botella actuales, se diseña y se construye un modelo de información real y funcional. Es decir, los conceptos informáticos se aplican para crear una herramienta tecnológica que mejore la gestión diaria en Bellidental.

El diseño del estudio es preexperimental porque trabaja con un único grupo de análisis al cual se le aplica una medición antes y después de usar el nuevo sistema. En este caso, el grupo está conformado por el personal administrativo y médico de la clínica. Primero se evalúa el tiempo de gestión actual de los pacientes usando los métodos tradicionales. Luego de implementar el software desarrollado, se vuelve a medir ese mismo tiempo operativo. Esta comparación directa es clave para comprobar el impacto real de la herramienta.

En conclusión, la combinación de estas tres modalidades metodológicas asegura que el proyecto tenga un respaldo científico sólido y resultados comprobables. El alcance aplicativo permite construir el sistema integrado que la clínica necesita, el diseño preexperimental proporciona la estructura para probarlo en el entorno de trabajo y el enfoque cuantitativo entrega los datos exactos para confirmar que el tiempo administrativo por paciente disminuyó de forma exitosa.